%! The Language of Variation — Unit 1, Lesson 1.1 — Slides (Beamer)
\documentclass[aspectratio=169, 11pt]{beamer}
\usepackage{apstats-beamer}   % provides \navyheader{} and \sectionlabel[color]{}
\usetikzlibrary{positioning}  % -beamer loads tikz but not this library

% The C and Y column types live in apstats-article, which a beamer deck does not
% load — redeclare them here rather than pulling in the article preamble.
\newcolumntype{C}[1]{>{\centering\arraybackslash}p{#1}}
\newcolumntype{Y}{>{\centering\arraybackslash}X}

\begin{document}

% TITLE SLIDE (CourseName is not defined in beamer, so the name is literal)
{
\setbeamercolor{background canvas}{bg=navy}
\begin{frame}[plain]
  \vspace{2.8cm}\hspace{0.55cm}%
  \begin{minipage}{0.9\textwidth}
    {\color{goldacc}\bfseries\small LESSON 1.1}\\[0.5cm]
    {\color{white}\bfseries\LARGE The Language of Variation}\\[1.2cm]
    {\color{skymid}\small Statistics~$\cdot$~Unit 1}
  \end{minipage}
\end{frame}
}

% The deck follows the EFFL flow: targets (SPOILER-FREE) → warm-up → activity launch →
% debrief frames (formalize in "red ink") → QuickNotes arc → application → practice →
% homework. Everything before the debrief divider stays vocabulary-free: the words
% `categorical' and `quantitative' may not appear until after the divider.
\newcommand{\redink}[1]{{\color{redacc}\bfseries #1}}

% ── LEARNING TARGETS (spoiler-free: no `categorical', no `quantitative') ─────
\begin{frame}[t, plain]
  \navyheader{Today: experience first, formalize later}
  \sectionlabel{Learning targets}
  By the end of today, I can\ldots
  \begin{itemize}\setlength{\itemsep}{5pt}
    \item sort the columns of a table into two kinds that behave differently.
    \item explain why adding up one column is useful and adding up another is nonsense.
    \item say what a column really records, even when every entry is written in digits.
    \item say what happens to a column when amounts get replaced by group names.
  \end{itemize}
  \begin{block}{How today works}
    Warm-up $\to$ \textbf{the trailhead register} with your group $\to$ we name what you
    found $\to$ practice. You will figure it out \emph{before} I give you the words.
  \end{block}
\end{frame}

% ── HOOK ─────────────────────────────────────────────────────────────────────
{
\setbeamercolor{background canvas}{bg=navy}
\begin{frame}[plain]
  \vspace{1.4cm}
  \begin{center}
    {\color{skymid}\bfseries\small TWO JERSEY NUMBERS: 12 AND 88}\\[0.9cm]
    {\color{white}\bfseries\fontsize{60}{60}\selectfont 50}\\[0.8cm]
    {\color{goldacc}\small That is their average. What did we just learn about those
    two players?}
  \end{center}
\end{frame}
}

% ── ACTIVITY LAUNCH ──────────────────────────────────────────────────────────
\begin{frame}[t, plain]
  \navyheader{Experience \& Formalize: The Trailhead Register}
  \sectionlabel[goldacc]{Groups of 3--4 $\cdot$ 20 minutes}
  Every party that hikes at Ross Hollow signs in and signs out. The ranger keeps the sheet
  all season.
  \begin{itemize}\setlength{\itemsep}{5pt}
    \item \textbf{Problem 1} --- sort the columns into two piles \emph{of your own}, and
          name them yourselves.
    \item \textbf{Problem 2} --- four more columns, every one written in digits.
  \end{itemize}
  \begin{block}{Your job}
    Use only what you already know. For \emph{every} answer, be ready to show me
    \textbf{where in the register} you see it.
  \end{block}
\end{frame}

% ── DEBRIEF DIVIDER — everything after this point may name the vocabulary ─────
{
\setbeamercolor{background canvas}{bg=navy}
\begin{frame}[plain]
  \vspace{1.8cm}\hspace{0.8cm}%
  \begin{minipage}{0.86\textwidth}
    {\color{goldacc}\bfseries\small NOW WE NAME IT}\\[0.7cm]
    {\color{white}\bfseries\Large You already built the two piles. We are only giving them
    the names statisticians use.}\\[0.9cm]
    {\color{skymid}\small Fill in your QuickNotes box as each name lands.}
  \end{minipage}
\end{frame}
}

\begin{frame}[t, plain]
  \navyheader{Two kinds of variable}
  \sectionlabel{your two piles, named}
  \vspace{2pt}
  \begin{columns}[t]
    \begin{column}{0.47\textwidth}
      \begin{block}{Categorical variable}
        Values are \textbf{category names} or \textbf{group labels}.
      \end{block}
      \vspace{4pt}
      {\small Trail $\cdot$ Dog along $\cdot$ Weather\\[3pt]
      \emph{Which group is this party in?}}
    \end{column}
    \begin{column}{0.47\textwidth}
      \begin{block}{Quantitative variable}
        Values are \textbf{numerical}, for a \textbf{measured or counted} quantity.
      \end{block}
      \vspace{4pt}
      {\small Group size $\cdot$ Time out (min)\\[3pt]
      \emph{How much? How many?}}
    \end{column}
  \end{columns}
  \vspace{12pt}
  \begin{center}
    {\small\color{slate} A \textbf{variable} is a characteristic that changes from one
    individual to another. Every variable is one kind or the other.}
  \end{center}
\end{frame}

% ── THE CRUX ─────────────────────────────────────────────────────────────────
\begin{frame}[t, plain]
  \navyheader{Digits are not the test}
  \sectionlabel{the error this lesson exists to kill}
  \vspace{4pt}
  \renewcommand{\arraystretch}{1.5}
  \begin{tabularx}{\textwidth}{@{} l C{3.6cm} X @{}}
    \hline
    \textbf{Column} & \textbf{Average of 2041 \& 6318} & \textbf{What it means} \\
    \hline
    Permit \#      & 4179.5 & \redink{Nothing.} Nobody had ``4179.5 of permit.'' \\
    Water bottles  & 4      & Between them, 4 bottles each. A real amount. \\
    \hline
  \end{tabularx}
  \vspace{14pt}
  \begin{block}{The test}
    Not \emph{``does it have digits?''} but \emph{``is the value an \redink{amount} of
    something?''} If adding or averaging is meaningful, it is quantitative.
  \end{block}
\end{frame}

\begin{frame}[t, plain]
  \navyheader{Written in digits, still categorical}
  \sectionlabel{the whole family}
  \vspace{6pt}
  \begin{columns}[t]
    \begin{column}{0.48\textwidth}
      \begin{itemize}\setlength{\itemsep}{7pt}
        \item ZIP code
        \item Jersey number
        \item Student ID number
        \item Parking spot number
        \item Phone number
        \item Permit number
      \end{itemize}
    \end{column}
    \begin{column}{0.48\textwidth}
      \begin{block}{Essential Knowledge (VAR-1.C.1)}
        A categorical variable takes on values that are \textbf{category names or group
        labels}.
      \end{block}
      \vspace{6pt}
      {\small Each of these is a \textbf{name that happens to be made of digits}. Sorting
      them works; averaging them does not.}
    \end{column}
  \end{columns}
\end{frame}

% ── THE REVERSE TRAP ─────────────────────────────────────────────────────────
\begin{frame}[t, plain]
  \navyheader{``62\% of parties hiked with a dog''}
  \sectionlabel{a number --- but from which kind of column?}
  \vspace{8pt}
  \begin{itemize}\setlength{\itemsep}{9pt}
    \item Every entry in the \emph{Dog along} column is \textbf{Yes} or \textbf{No}.
    \item Nobody wrote a number in it. The 62\% is something \emph{we built} by counting.
    \item So \emph{Dog along} is still \redink{categorical}.
  \end{itemize}
  \vspace{10pt}
  \begin{block}{Remember this one}
    Counting a categorical variable \textbf{always} produces numbers. The \textbf{summary}
    is numeric; the \textbf{variable} is not.
  \end{block}
\end{frame}

% ── HOW EACH KIND IS SUMMARIZED ──────────────────────────────────────────────
\begin{frame}[t, plain]
  \navyheader{What you are allowed to do with each}
  \sectionlabel{this decides every graph for the rest of the course}
  \vspace{18pt}
  \begin{center}
  \begin{tikzpicture}[node distance=0.6cm, every node/.style={font=\small\bfseries}]
    \node[fill=navy, text=white, rounded corners=2mm, minimum height=1.1cm,
          minimum width=3.4cm] (v) {A variable};
    \node[fill=skymid, text=charcoal, rounded corners=2mm, minimum height=1.1cm,
          minimum width=3.4cm, above right=0.35cm and 1.6cm of v] (c) {Categorical};
    \node[fill=goldacc, text=charcoal, rounded corners=2mm, minimum height=1.1cm,
          minimum width=3.4cm, below right=0.35cm and 1.6cm of v] (q) {Quantitative};
    \node[fill=greenacc, text=white, rounded corners=2mm, minimum height=1.1cm,
          minimum width=4.6cm, right=of c] (cs) {Count $\to$ percents};
    \node[fill=greenacc, text=white, rounded corners=2mm, minimum height=1.1cm,
          minimum width=4.6cm, right=of q] (qs) {Average, spread};
    \draw[->, thick, navy] (v) -- (c);
    \draw[->, thick, navy] (v) -- (q);
    \draw[->, thick, navy] (c) -- (cs);
    \draw[->, thick, navy] (q) -- (qs);
  \end{tikzpicture}
  \end{center}
  \vspace{14pt}
  \begin{center}
    {\small Pick the wrong branch and every graph and statistic after it is wrong too.}
  \end{center}
\end{frame}

% ── BINNING ──────────────────────────────────────────────────────────────────
\begin{frame}[t, plain]
  \navyheader{One variable can be recorded either way}
  \sectionlabel{and the choice costs you something}
  \vspace{6pt}
  \begin{columns}[t]
    \begin{column}{0.46\textwidth}
      \begin{block}{Age: 14, 27, 41, 63}
        Quantitative. You can average it, and you know 41 is 14 years older than 27.
      \end{block}
    \end{column}
    \begin{column}{0.46\textwidth}
      \begin{block}{Age band: teen, adult, senior}
        Categorical. You can count each band --- and that is all.
      \end{block}
    \end{column}
  \end{columns}
  \vspace{14pt}
  \begin{center}
    {\large Same people. Same information. \redink{Different kind of variable.}}\\[8pt]
    {\small\color{slate} Grouping amounts into labels is a one-way door: you cannot
    get the mean back out.}
  \end{center}
\end{frame}

% ── C or Q? ──────────────────────────────────────────────────────────────────
\begin{frame}[t, plain]
  \navyheader{C or Q?}
  \sectionlabel{call it out}
  \vspace{4pt}
  \renewcommand{\arraystretch}{1.7}
  \begin{tabularx}{\textwidth}{@{} X C{3.2cm} @{}}
    \hline
    \textbf{Variable} & \textbf{C or Q} \\
    \hline
    Number of siblings a student has & \\
    The student's home ZIP code      & \\
    Shirt size (S / M / L)           & \\
    Time to run 400 m, in seconds    & \\
    \hline
  \end{tabularx}
  \vspace{10pt}
  \begin{center}
    {\small\color{slate} For every one you call \textbf{Q}: what would its average
    actually mean?}
  \end{center}
\end{frame}

% ── APPLICATION ──────────────────────────────────────────────────────────────
\begin{frame}[t, plain]
  \navyheader{Application: The Season Pass}
  \sectionlabel[goldacc]{We work this one together --- you hold the pen}
  A ski area records, for each of 3{,}400 season-pass holders: \textbf{pass number},
  \textbf{age}, \textbf{pass tier} (Gold / Silver / Bronze), \textbf{miles traveled}, and
  \textbf{locker number}.
  \begin{itemize}\setlength{\itemsep}{5pt}
    \item Label all five: C or Q.
    \item Age in years becomes \emph{age group}. Which kind now --- and what did we lose?
    \item ``The average pass number is 5{,}120.'' Why does that tell us nothing?
  \end{itemize}
  \begin{block}{Ask it every time}
    Is this value an \redink{amount of something}, or is it the \redink{name of a group}?
  \end{block}
\end{frame}

% ── PRACTICE (in class, not scored) ──────────────────────────────────────────
\begin{frame}[t, plain]
  \navyheader{Check Your Understanding}
  \sectionlabel{5 exercises $\cdot$ pairs first, then on your own}
  \begin{itemize}\setlength{\itemsep}{5pt}
    \item A coffee shop's order log: classify four columns.
    \item Number of pets vs.\ type of pet --- and how each is summarized.
    \item Body temperature vs.\ Normal/Fever: where did the second column come from?
    \item Name a variable in digits that is not quantitative.
    \item One multiple-choice item, with a one-sentence justification.
  \end{itemize}
  \begin{block}{This is practice, not a quiz}
    These are \textbf{not scored}. Start with 1, 3, and 5. I will look over your shoulder
    at 5 on the way out.
  \end{block}
\end{frame}

% ── HOMEWORK ─────────────────────────────────────────────────────────────────
\begin{frame}[t, plain]
  \navyheader{AP Practice, homework \& what's next}
  \sectionlabel[goldacc]{AP Practice --- practice, not a grade}
  \begin{itemize}\setlength{\itemsep}{5pt}
    \item Four multiple-choice items on the trailhead register.
    \item One free-response set: a city library's 900 cardholders.
    \item Optional: bring in any data table and classify every column.
  \end{itemize}
  \begin{block}{Your homework is in DeltaMath}
    The \redink{graded} assignment is posted in \textbf{DeltaMath} --- log in and complete it
    before next class. The AP Practice page is not scored: work it, and bring what you get
    stuck on to study hall or office hours.
  \end{block}
\end{frame}

% ── CLOSE ────────────────────────────────────────────────────────────────────
{
\setbeamercolor{background canvas}{bg=navy}
\begin{frame}[plain]
  \vspace{1.5cm}\hspace{0.8cm}%
  \begin{minipage}{0.86\textwidth}
    {\color{goldacc}\bfseries\small BEFORE YOU GO}\\[0.7cm]
    {\color{white}\bfseries\Large Digits do not make a variable quantitative. Ask whether
    the value is an amount of something.}\\[0.9cm]
    {\color{skymid}\small Next: Lesson 1.2 --- we take a categorical column and count it
    properly, with frequency and relative frequency tables.}
  \end{minipage}
\end{frame}
}

\end{document}
